\section{Resultados}

El juego final cumple todos los requisitos de la asignatura (Tabla~\ref{tab:requisitos}) y se
entrega en las tres plataformas pedidas.

\begin{TablaAPA}{tab:requisitos}{Cumplimiento de los requisitos de la asignatura}
\begin{tabularx}{\linewidth}{P{4.6cm}L}
\toprule
Requisito & Cómo lo cumple Nexus \\
\midrule
Tres niveles jugables con progresión lógica & Cuarto de Mantenimiento, Pasillos de Almacenamiento y Núcleo de Datos; cada uno se desbloquea al extraer el disco del anterior \\
Tres Canvas diferenciados & Menú principal, HUD y pantalla de Game Over y Victoria, además del inventario \\
Selector de personajes con Avatar Humanoid & Avatares de los dos integrantes configurados como Humanoid, con un único controlador de animación \\
Música de fondo con componentes de audio & Música por capas (menú, exploración, tensión, jefe, alarma y victoria) y ambiente del centro de datos \\
Créditos e imagen institucional & Pantalla de créditos con la materia, la Universidad Central, el docente y los integrantes; logo en el título y en los créditos \\
Explicación del objetivo & Informe de la misión antes del Nivel~1 y objetivo actual en el HUD \\
Temporizador & Refrigeración con cuenta atrás y temperatura del núcleo \\
Sistema de vidas o salud & Tres oportunidades de evasión con estado herido y objetos de curación \\
Objetos coleccionables con contador & Discos Duros Maestros (x/3) y cartuchos de refrigerante \\
Inteligencia artificial autónoma & Androides con patrulla, visión, oído, persecución y búsqueda sobre una malla de navegación, y un jefe final \\
Cuadros de diálogo narrativos & Terminales, notas, avisos del sistema y subtítulos de las cinemáticas \\
Plataformas & Windows (\cod{.exe}), WebGL con pantalla completa y Android (\cod{.apk}) \\
\bottomrule
\end{tabularx}
\end{TablaAPA}

Además de los requisitos, el juego incorpora una habilidad distinta para cada personaje, armas con
munición limitada y apuntado asistido, cadáveres que se registran con notas que cuentan la historia,
animaciones de escopeta creadas por código, cámaras que guían al jugador hacia los objetos
importantes, un inventario con archivo y mapa, guardado en terminales, máquinas expendedoras con
riesgo, un puzzle de válvulas, un jefe final, cinemáticas hechas con el motor, una transición
animada entre niveles y controles táctiles.

\section{Conclusiones}

El desarrollo por fases, con una versión jugable al final de cada una y compilaciones de prueba
desde el inicio, permitió construir un juego completo de tres niveles y llevarlo a las tres
plataformas sin sorpresas al final del proyecto.

La construcción del contenido mediante herramientas de editor programadas fue la decisión que más
influyó en el resultado. Describir los niveles como mapas de texto y generar texturas, objetos,
decoración, cámaras y cinemáticas por código hizo posible producir escenarios grandes y coherentes
con un equipo de dos personas, y aplicar cualquier ajuste de diseño en todo el juego con solo
reconstruir.

El sistema de Avatar Humanoid de Unity fue la base del trabajo de personajes: permitió que los
avatares de los integrantes, los androides y los personajes secundarios compartieran animaciones,
y que el selector de personajes funcionara con un único controlador. Ese mismo sistema hizo posible
crear poses y animaciones nuevas músculo por músculo, como las de los cadáveres y la escopeta.

Las referencias del survival horror clásico orientaron las decisiones de diseño: cámaras fijas que
cortan cuando el personaje sale del cuadro, controles tipo tanque, recursos escasos, sigilo frente
a enemigos que ven y oyen, y una mezcla de audio que acompaña la tensión de cada situación.

Finalmente, adaptar el juego a celulares y navegadores exigió cuidar el rendimiento desde el
principio (mallas combinadas, sombreadores sin scripts, una sola cámara activa) y añadir controles
táctiles, lo que dejó un juego que se puede jugar de la misma forma en las tres plataformas.
